\section{复议}
\begin{enumerate}
  \item \zd{复议本质上是个具体行政行为}。95\% 和诉讼一样。
  \item 被告叫\zd{复议被申请人}。\bs{做题/选项要看清问/说的是哪个主体}。
  \item 居中裁判的叫\zd{复议机关}。\bs{原则找上级}。找爹。
        \begin{enumerate}
          \item 1 个。
                \begin{enumerate}
                  \item 政府 \yjt 上级政府。例外。\zd{省部级原级复议}。
                  \item 国家垂直领导机关。国安；\zd{海关、税务}（收钱）；金融、外汇（花钱）。
                  \item \zd{派出所} \yjt 爹 \hz 爷爷。公安局和政府。
                \end{enumerate}
          \item 2个。政府的部门 \yjt \zd{业务上级 + 爹}。\al{县工商局 \yjt 市工商局、县政府}。
          \item \zd{经批准行为复议看上级}。
        \end{enumerate}
  \item \zd{仅公民可以口头委托} 1 - 2 个代理人。\al{企业 \cuo}。
  \item 告漏了。直接\zd{强行追加}。
  \item 程序。\zd{时间注意不要记混，复议不能小于规定时间，诉讼可以}。
        \begin{enumerate}
          \item \bs{起诉 60 受理 5}。起诉期 \dydy 60 日（只多不少）。受理 5 日。不受理 \yjt 继续找上级（\zd{催 \hz 越级受理}）。
          \item \bs{决定 60 举证 10}。决定 \xydy 60 日（高效便民）。举证期 10 日。
          \item \zd{2} 人审理。适用\zd{实质审查}。
          \item 审理方式。\yi 原则书面。\er \zd{重大复杂听证}。主动 \hz \zd{依原告申请}。\al{以开庭方式审理、被申请人申请听证 \cuo}。
        \end{enumerate}
  \item 特别程序。
        \begin{enumerate}
          \item 中止暂时停。终止\zd{曲终人散}。
                \begin{enumerate}
                  \item \bs{人死看继承。60 中转终，诉讼等 90}。
                  \item \bs{复议行转刑是终止，诉讼是中国}（止）。
                \end{enumerate}
          \item 和解、调解。同诉讼。\zd{可强制执行}。
          \item 撤回。再次复议不受理，但可去诉讼。例外。\al{胁迫}。
          \item 改变原行政行为。\zd{只要不撤诉照审不误}。
        \end{enumerate}
  \item 判决。送达生效。\bs{直接说怎么判的大概率都不对，因为要先判断合法、不合法}。
        \begin{enumerate}
          \item 合法。维持决定、\zd{不作为}驳回。\al{罚款、不发证}。
          \item 不合法。事实、程序、法律、内容\bs{所有的行为都可撤可变}。仅证据问题需要先查清，再处理。
          \item 抽象行为附带审查。\bs{能管就管，不能就转}。
          \item 被申请人\zd{不履行、拖延履行}复议决定，复议机关 \hz 上级应当\zd{责令其限期履行}。
        \end{enumerate}
  \item 复议期间发现\zd{被申请人}的相关行政行为违法。\bs{个别}/普遍现象 \yjt 可以制作\bs{意见书}/建议书。
  \item \zd{当事人可自行委托鉴定，费用自担}。
\end{enumerate}
